\section{The eighteen local inputs to the automorphy argument}

% BEGIN V72L1
\paragraph{V72L1. Finiteness and the compact pairing.}
Write $G_n=\mathrm{GL}_n(k)$ for a nonarchimedean local field $k$,
$P=MU$ for an upper block parabolic, $i_P$ for parabolic induction
with the factor $\delta_P^{1/2}$, and $r_P$ for unipotent
coinvariants twisted by $\delta_P^{-1/2}$. All representations here
are smooth over $\mathbb C$. An irreducible representation of $G_n$
is admissible. Both $i_P$ and $r_P$ preserve finite length.
For $i_P$ this follows first on irreducibles from finiteness of
the length of parabolic induction, then on a composition series by
exactness. For $r_P$ one can embed an irreducible in an induction
from a cuspidal representation and use the finite Bruhat filtration
with cuspidal input. Applying exactness again treats all finite
lengths. These are the finiteness results of Bernstein and Zelevinsky,
\emph{Induced representations of reductive $p$-adic groups I}
(1977), \S\S2 and 5.

The fourth assertion in \texttt{bz} specifies the actual smooth
duality isomorphism for an arbitrary smooth $M$-module $W$. If
$K=\mathrm{GL}_n(\mathcal O_k)$, its pairing is
\[
 (f',f)\longmapsto\int_K f'(x)(f(x))\,dx,
 \qquad f'\in i_P(W^\vee),\quad f\in i_P(W).
\]
The two half-modulus factors give the transformation by $\delta_P$
required by integration on $P\backslash G_n$. To see that this gives
the whole smooth dual, fix a compact open subgroup $J\subset K$.
In the compact realization $G_n=PK$, a $J$-invariant function is a
tuple over the finite set $(P\cap K)\backslash K/J$. The entry at $x$
belongs to the fixed space of the image in $M$ of
$P\cap xJx^{-1}$. Averaging over this compact open image extends
every linear form on its fixed space to an invariant smooth linear
form on $W$. On these tuples the integral is the direct sum of
evaluation pairings with positive coset volumes. Every smooth dual
vector is fixed by some such $J$, proving the asserted isomorphism
for every Haar normalization. Zero-sized blocks contribute trivial
groups, including $G_0$.

Bernstein's unitary induction theorem supplies \texttt{bern}: the
external tensor product of two irreducible unitarisable factors is
an irreducible unitary Levi representation, and its normalized
induction is irreducible. The invariant positive forms give the
unitary completions, whose smooth vectors recover the admissible
representations. A rank-zero factor is the one-dimensional unit.
See Bernstein, \emph{$P$-invariant distributions on $\mathrm{GL}(N)$
and the classification of unitary representations}, 1984, p.~51.
% END V72L1

% BEGIN V72L2
\paragraph{V72L2. The geometric lemma and its multiplicities.}
Take parabolics with block sizes $(a_i)$ and $(b_j)$. The relative
positions of their flags are precisely the matrices $C=(c_{ij})$ of
nonnegative integers with those row and column sums: $c_{ij}$ is the
dimension of the successive quotient of the $(i,j)$ intersection
of the two flags. Order the orbits so that successive unions used
for the support filtration are open in the part remaining. The
geometric lemma of Bernstein and Zelevinsky (I, 2.12 and \S5)
gives a filtration of $r_Qi_P(W)$ with one quotient per matrix.
For this quotient take Jacquet modules in every row with blocks
$c_{ij}$, transpose the resulting array of groups, then induce in
each column. Restriction to the orbit followed by taking unipotent
coinvariants gives this quotient; the modular characters of the
stabilizers are exactly those in normalized restriction and
induction. The filtration order depends on the flags and works for
every smooth $W$. Changing coordinates on a row or column conjugates
the intermediate parabolic and gives an isomorphic quotient.

For $W=\boxtimes_i W_i$ of finite length let $T$ consist of the
arrays of irreducibles $s_{ij}$ for which every
$\boxtimes_j s_{ij}$ occurs in $r_{C_i}W_i$. This is a finite set
by V72L1 and factorization of irreducible admissible modules of
a product as external tensor products. Off $T$ some row multiplicity
vanishes. On $T$, multiply the row multiplicities, then multiply by
the multiplicity of the desired constituent in the external tensor
product of the column inductions. Summing over $T$ gives the
multiplicity in the cell. Indeed tensor product over $\mathbb C$
and normalized induction are exact, so one can apply them successively
to composition series. This proves both clauses of \texttt{geo}.
% END V72L2

% BEGIN V72L3
\paragraph{V72L3. Two conventions for segments and their products.}
Let $\nu=|\det|_k$ and
$\Delta=[\rho,\nu\rho,\ldots,\nu^{l-1}\rho]$, with $\rho$
supercuspidal of degree $d>0$ and $l>0$. In this ascending product
the irreducible submodule is $Z(\Delta)$ and the irreducible
quotient is $L(\Delta)$. In particular these two symbols have
different meanings even for a two-term segment. The normalized
Jacquet module of $Z(\Delta)$ for the $l$ blocks of size $d$ is
the single ordered tensor of the terms of $\Delta$. Conversely that
Jacquet module yields, by Frobenius reciprocity, a nonzero map of
the irreducible into the ascending product, identifying it with
$Z(\Delta)$. The ring of finite-length classes, with normalized
induction as product, is the integral polynomial ring on the
$[Z(\Delta)]$. Its underlying additive group has the irreducible
classes as basis by Jordan--H\"older; the unit is the rank-zero
class. Sending each $[Z(\Delta)]$ to its indeterminate gives
exactly \texttt{zel}. See Zelevinsky, \emph{Induced representations
of reductive $p$-adic groups II} (1980), \S3 and Corollary 7.5.

For \texttt{disj}, write the two irreducibles in the $Z$ classification.
Their cuspidal supports are the multisets obtained by reading all
entries in their segments. If a segment on each side formed a linked
pair, they would overlap or abut on the same cuspidal line. This gives
equal entries or entries differing by $\nu^{\pm1}$, in the same degree,
contrary to the stated hypothesis. Zelevinsky II, Proposition 8.5,
therefore gives irreducibility of the product.

For the quotient convention, put the segments in an order in which
no earlier segment precedes a later one. A linked segment precedes
another when the latter starts strictly to its right on the same
cuspidal line. The ordered induction of their $L(\Delta)$ has a
simple head, hence a unique maximal submodule (\texttt{cos}).
The head $L(\mathfrak m)$ labels every irreducible exactly once
(\texttt{exh} and the first clause of \texttt{multi}). This is the
Langlands classification for $G_n$: ordering the essentially
square-integrable factors by decreasing real central exponent gives
a standard module. Any order satisfying the given precedence
condition can be reached from this order by swapping adjacent
unlinked factors; their intertwining maps are isomorphisms.
The empty multisegment yields the rank-zero unit. Finally
$L(\mathfrak m+\mathfrak n)$ occurs in
$L(\mathfrak m)\times L(\mathfrak n)$. In the $Z$ convention the
sum occurs with multiplicity one (Zelevinsky II, Proposition 8.4).
The Zelevinsky involution interchanges the $Z$ and $L$ bases and is
multiplicative on the Grothendieck ring, giving the claimed occurrence
in the $L$ convention. This uses the involution theorem of
M\oe glin and Waldspurger, \emph{Sur l'involution de Zelevinski}
(1986). A positive composition multiplicity gives precisely the
subquotient in the second clause of \texttt{multi}.
% END V72L3

% BEGIN V72L4
\paragraph{V72L4. Square-integrability, temperance and genericity.}
For an admissible irreducible with central character $\omega$, the
absolute values of its smooth matrix coefficients descend to $G_n/Z$
when $|\omega|=1$. Conversely the descent required in
\texttt{CoeffLp} forces $|\omega|=1$: a coefficient nonzero at the
identity takes value $\omega(z)$ times that value at $z$. Smooth
linear forms separate vectors by compact open averaging.
Square-integrability is then the condition that these functions
belong to $L^2(G_n/Z)$, and the tempered criterion is membership
in every $L^{2+\varepsilon}$ for $\varepsilon>0$. Any two quotient
Haar measures are positive scalar multiples, so the choice of
measure does not change these conditions.

The essentially square-integrable irreducibles of positive degree
are the $L(\Delta)$ (Zelevinsky II, Theorem 9.3).
A determinant twist by a real power of the normalized absolute
value cancels the real central exponent and gives a square-integrable
representation. Conversely, twisting a square-integrable irreducible
back by a smooth determinant character gives an essentially
square-integrable irreducible. This proves \texttt{seg}. The
unitary-dual definition in \texttt{sq} requires positive rank,
unitarisability and square-integrable coefficients. For such a
representation $|\omega|=1$, so evaluating the coefficient at the
chosen representative of a central coset gives the same absolute
value at every representative. It satisfies the preceding criterion
with the trivial twist and hence is an $L(\Delta)$.

The tempered classification and the matrix-coefficient criterion
(Casselman, \emph{Introduction to the theory of admissible
representations of $p$-adic reductive groups}, \S\S4--6) realize
each tempered irreducible as a constituent of an induction of
unitary square-integrable representations of a Levi. The induction
is irreducible by V72L1, applied successively, and normalized
induction of tempered representations is tempered. This gives the
equivalence \texttt{temp}, with positive block sizes summing to $n$.
Finally a generic irreducible is itself a product of $L(\Delta)$
for pairwise unlinked segments, by Zelevinsky II, Theorem 9.7(b).
The source uses a nontrivial continuous additive character and a
nonzero Whittaker functional. These are precisely its generic
representations. The rank-zero case uses the empty product.
This supplies \texttt{gen}.
% END V72L4

% BEGIN V72L5
\paragraph{V72L5. The canonical map at depth zero.}
Let $K^1$ be the kernel of reduction on $K$, let $U$ be the block
unipotent radical, and set $J=K^1(U\cap K)$. Reduction identifies
$J$ with the inverse image of the unipotent radical in the finite
parabolic: lift the upper blocks entry by entry. Thus $J$ is the
pro-unipotent radical of the lattice-chain parahoric attached to
the block flag, and $J\cap M$ is the product of the principal
congruence subgroups of the blocks. Proposition 6.7 of Moy and
Prasad, \emph{Jacquet functors and unrefined minimal $K$-types}
(1996), applies to this facet and its associated Levi. For every
admissible representation $V$ its canonical map
\[
                  V^J\longrightarrow(V_U)^{J\cap M}
\]
is an isomorphism. This is the quotient map on vectors; it gives
the specific bijection in \texttt{mp}. For a single nonempty block
the map is the identity. Permuting coordinates transports the
calculation to every indexing equivalence in the statement.

For \texttt{dzi}, a cuspidal support all of whose factors have
principal-congruence fixed vectors has depth zero in the product
Levi. By Moy--Prasad, Theorem 5.2(1), each irreducible subquotient
of its parabolic induction has the same depth. A nonzero vector of
depth zero is fixed by the pro-unipotent radical attached to some
facet. Choosing a vertex in its closure gives a principal congruence
subgroup contained in that radical. Indeed for a lattice chain,
the kernel at the vertex acts trivially on every successive quotient
of the chain. Every vertex of the reduced building of $G_n$ is
represented by a lattice, and $G_n$ carries it to the standard
lattice. Transporting the vector gives a nonzero $K^1$-fixed vector.
The empty-rank representation satisfies the conclusion directly.
% END V72L5

% BEGIN V72L6
\paragraph{V72L6. The rank-difference-one converse.}
Put $n=m+1$, $m>0$, and fix the indicated nontrivial additive
character. The hypothesis gives the same scalar $e_\tau$ for two
cuspidal representations against each irreducible of $G_m$. The
literal functional equation quantifies over all invariant integrals
on test functions. The following argument from the source shows
how that specific hypothesis determines the representation.

Choose nonzero Whittaker functionals on the two cuspidals, writing
$W_i$ for their translates. Write $c(g)=\operatorname{diag}(g,1)$
and $\widetilde g=w\,{}^tg^{-1}$ with the long Weyl element of the
relevant rank. Suppose $W_1(c(g))=W_2(c(g))$ for all $g$. The
difference $D(g)=W_1(\widetilde{c(g)})-W_2(\widetilde{c(g)})$
is a locally constant Whittaker test function of compact support
modulo the maximal unipotent subgroup, by cuspidality. Use the
Whittaker separation lemma: on these test functions the invariant
pairings with opposite-character Whittaker functions of irreducible
$G_m$-modules separate points. Choose its integral $I$ in the
hypothesis. To test against a translate of a functional on $\sigma$,
take $\tau$ to be $\sigma$ pulled back by inverse transpose with
the source's Weyl conjugation. Translating the test vector by the
inverse image of the long element turns its tilded test function
into the original translate for $\sigma$. Subtract the two functional
equations. The right sides coincide because the corner functions,
the central character of $\tau$ and $e_\tau$ coincide. Hence all
separating pairings of $D$ vanish. In particular $D(1)=0$, giving
equality at $w_n$ for every pair agreeing on the corner.

There exists a pair agreeing on the corner and nonzero at the
identity: normalize the two functionals on chosen vectors to take
value one, then apply the Kirillov restriction lemma to their
direct sum, using the difference functional. The big-cell argument
now gives equality of the central characters and propagates agreement
on the corner to agreement on $G_n$. The pairs with equal Whittaker
functions form an invariant subspace of $V_1\oplus V_2$. Its two
projections are nonzero and hence surjective by irreducibility.
Their kernels vanish because the translates of each nonzero
Whittaker functional separate vectors in an irreducible representation.
The two projections thus identify the representations. These are the
separation, Kirillov and big-cell steps assembled in
\texttt{LConvFinal}; they establish \texttt{conv} with the single
scalar and the test quantifiers actually used here.
% END V72L6

% BEGIN V72L7
\paragraph{V72L7. Three character calculations.}
For \texttt{brauer}, let $A(G)$ denote the subgroup of the complex
character ring generated over $\mathbb Z$ by
$\operatorname{Ind}_H^G(\lambda-1)$ for linear characters $\lambda$
of subgroups $H$. For an elementary group $D$ every irreducible is
monomial. If $\alpha=\operatorname{Ind}_S^D\lambda$, then
$\alpha-\alpha(1)$ differs from
$\operatorname{Ind}_S^D(\lambda-1)$ by
$\operatorname{Ind}_S^D1-[D:S]$. The latter belongs to $A(D)$:
choose a chain from $S$ to $D$ with successive normal inclusions of
prime index, possible because $D$ is nilpotent. At each step the
permutation character is the inflated regular character of a cyclic
group, and subtracting its degree gives a sum of linear characters
minus one. Induction in stages proves the assertion along the chain,
using
\[
 \operatorname{Ind}_S^D1-[D:S]
 =\operatorname{Ind}_H^D(\operatorname{Ind}_S^H1-[H:S])
   +[H:S](\operatorname{Ind}_H^D1-[D:H]).
\]
Thus every virtual degree-zero character of $D$ belongs to $A(D)$.
Apply Brauer's induction theorem to the trivial character of $G$,
writing it as an integral combination of inductions of linear
characters from elementary groups. Multiply by
$\chi-\chi(1)$ and use the induction tensor identity. Each term is
induced from a virtual character of degree zero on an elementary
group, proving $\chi-\chi(1)\in A(G)$. Evaluation at $g$ uses the
ordinary induced character sum with divisor $|H|$, exactly
\texttt{indFun}.

For \texttt{chev}, apply Chevalley's first theorem in
\emph{Deux théorèmes d'arithmétique} (1951), pp.~36--37, to the
finitely generated multiplicative group $\mathcal O_E^\times$,
the integer $m>0$, and the specified integer $b>0$. It supplies a
rational integer prime to $b$ for which units multiplicatively
congruent to $1$ are $m$-th powers in this same group. Its proof
constructs a nonzero modulus, whose absolute value gives $a>0$.
Divisibility of $u-1$ by $a$ in $\mathcal O_E$ implies the required
multiplicative congruence, with denominator one.

For \texttt{green}, let $q$ be the stated prime and $K$ the given
degree-$d$ field over $\mathbb F_q$, $d>0$. Its multiplicative group
is the elliptic torus in $\mathrm{GL}_d(\mathbb F_q)$. The Green
parametrization associates cuspidal irreducible characters bijectively
to Frobenius orbits of characters $\theta$ of $K^\times$ of orbit
length $d$. The degree is $\prod_{i=1}^{d-1}(q^i-1)$. An element
with irreducible characteristic polynomial has as its eigenvalues the
conjugates $z,z^q,\ldots,z^{q^{d-1}}$ of any one root. The character
there equals
\[
              (-1)^{d-1}\sum_{i=0}^{d-1}\theta(z^{q^i}).
\]
Changing the root permutes the summands. Cuspidality in the source
is irreducibility and vanishing of the dimensions of invariants
along the unipotent radicals of proper maximal parabolics, the
ordinary cuspidality condition. Set the unrestricted value of the
parameter map on nonregular characters equal to zero. This constructs
the required total function, with precisely the required equality
and exhaustion on regular characters. See Green, \emph{The characters
of the finite general linear groups}, Trans. Amer. Math. Soc.
\textbf{80} (1955), 402--447.
% END V72L7

% BEGIN V72L8
\paragraph{V72L8. Supplying the local record.}
Assign the fields of \texttt{LitLocAlt} as follows. V72L1 supplies
\texttt{bz} and \texttt{bern}, V72L2 supplies \texttt{geo},
V72L3 supplies \texttt{zel}, \texttt{disj}, \texttt{multi},
\texttt{exh} and \texttt{cos}, V72L4 supplies \texttt{temp},
\texttt{seg}, \texttt{sq} and \texttt{gen}, V72L5 supplies
\texttt{mp} and \texttt{dzi}, V72L6 supplies \texttt{conv}, and
V72L7 supplies \texttt{brauer}, \texttt{chev} and \texttt{green}.
This gives the local field of the literal automorphy package.

The source groups these proofs in a constructor taking
\texttt{OpenImageCondition}, whose only field is the global
Galois representation statement \texttt{RowT}. The local
constructor can be read independently. Its polynomial-ring wrapper
uses the closed segment construction, the ring theorem for those
segments, and finiteness. Its depth-zero induction wrapper calls the
geometric lemma, finiteness and that polynomial-ring result.
Its tempered wrapper uses the geometric lemma, the two finite-length
and admissibility results, the three Langlands classification fields
and the square-integrable segment theorem. Its unitary square-integrable
wrapper uses the latter theorem as well. Expanding these applications
leaves local theorems with no use of the global argument. The global
representation field of the separate automorphic record is still a
distinct assertion.
% END V72L8
